\section{ControlNet: ajuste fino de control condicional eficiente}

\subsection{Planteamiento del problema: lograr nuevas condiciones mediante modelos preentrenados}

Cuando se dispone de un modelo difusivo preentrenado en miles de millones de imágenes (como Stable Diffusion) y se desea que acepte nuevas entradas condicionales (por ejemplo, mapas de bordes, poses corporales o mapas de profundidad), surge una dificultad central:

\begin{warningbox}{El costo de entrenar desde cero un modelo condicional}
Si se entrena un nuevo modelo condicional desde cero utilizando \textbf{guía sin clase}:
\begin{itemize}
    \item Se requiere \textbf{datos emparejados} a escala de miles de millones (imagen condicional + imagen real), con un costo de recolección extremadamente alto
    \item Se necesita entrenar el modelo difusivo completo desde cero, lo que implica un costo computacional enorme
    \item No es posible reutilizar el amplio conocimiento visual aprendido por el modelo preentrenado existente
\end{itemize}
Por ejemplo, para entrenar un modelo de «mapa de bordes $\to$ imagen real», teóricamente se necesitarían 5.000 millones de pares (mapa de bordes, imagen real) de datos emparejados; en la práctica, es casi imposible obtener conjuntos de este tamaño y alta calidad.
\end{warningbox}

El objetivo de ControlNet es: \textbf{congelar los parámetros del modelo preentrenado y permitir que el modelo acepte nuevas entradas condicionales utilizando solo una cantidad reducida de datos emparejados (por ejemplo, entre 50K y 100K pares) y un número pequeño de parámetros adicionales}.

\subsection{Diseño arquitectónico de ControlNet}

El diseño de ControlNet se rige por dos principios fundamentales:

\textbf{Principio uno: mantener invariable la salida inicial (convolución cero)}

Al inicio del ajuste fino, la salida de los módulos añadidos debe ser cero, de modo que el comportamiento de toda la red sea idéntico al del modelo preentrenado. Esto se logra mediante \textbf{convolución cero}, es decir, inicializando tanto los pesos como los sesgos de las capas de conexión a cero.

\textbf{Principio dos: reutilizar el codificador preentrenado (copia entrenable)}

ControlNet crea una \textbf{copia entrenable} de la parte del codificador U-Net preentrenada, destinada a procesar las entradas condicionales:

\begin{enumerate}
    \item \textbf{U-Net original}: parámetros congelados, continúa procesando $z_t$ y $t$ (garantizando la calidad de generación)
    \item \textbf{Copia del codificador de ControlNet}: parámetros entrenables, recibe adicionalmente la imagen condicional $c$
    \item \textbf{Conexión por convolución cero}: la salida del codificador de la copia se suma a las conexiones salientes (skip connections) del U-Net original mediante una capa de convolución cero
\end{enumerate}

\begin{importantbox}{Estrategia de inicialización por convolución cero en ControlNet}
Al inicio del entrenamiento:
\begin{itemize}
    \item Todos los pesos y sesgos de la convolución cero son iguales a 0
    \item La salida de la copia de ControlNet queda «oculta» por la convolución cero, sin afectar al U-Net original
    \item La salida del sistema completo es \textbf{exactamente igual} a la del modelo preentrenado; las capacidades generativas ya aprendidas no se ven comprometidas por el añadido de nuevos módulos
\end{itemize}
A medida que avanza el entrenamiento, la convolución cero aprende pesos no nulos e incorpora gradualmente la información condicional al proceso de generación.
\end{importantbox}

\subsection{Entrenamiento de ControlNet}

Durante el entrenamiento solo se actualizan los parámetros de la copia del codificador de ControlNet y de las capas de convolución cero; el U-Net original permanece completamente congelado:

\begin{itemize}
    \item Datos de entrenamiento: conjunto reducido de datos emparejados (por ejemplo, 50K pares de mapa de bordes-imagen real)
    \item Objetivo de entrenamiento: pérdida difusiva estándar $\|\epsilon - \epsilon_\theta(z_t, t, c)\|^2$
    \item Costo de entrenamiento: considerablemente inferior al entrenamiento desde cero; generalmente se completan los entrenamientos en 8 GPUs durante varios días
\end{itemize}

\begin{knowledgebox}{Diversas modalidades condicionales en ControlNet}
Dado que el codificador condicional de ControlNet puede procesar cualquier entrada de imagen con la misma resolución que la imagen original, ha sido aplicado con éxito a:
\begin{itemize}
    \item \textbf{Mapa de bordes Canny}: controla contornos y bordes
    \item \textbf{Pose corporal (OpenPose)}: controla las acciones de las figuras humanas
    \item \textbf{Mapa de profundidad}: controla la estructura espacial y las relaciones de distancia
    \item \textbf{Segmentación semántica}: controla el contenido de las regiones
    \item \textbf{Bocetos/dibujos a mano alzada}: genera imágenes realistas a partir de trazos simples
    \item \textbf{Mapas normales}: controla la orientación de las superficies y la iluminación
\end{itemize}
Cada tipo de condición requiere entrenar un módulo de ControlNet específico, pero todos pueden compartir el mismo modelo base congelado.
\end{knowledgebox}

\subsection{Resumen de la sección}

Mediante tres diseños clave —congelación del modelo preentrenado, copia del codificador entrenable e inicialización por convolución cero—, ControlNet logra dotar a los modelos difusivos preentrenados de la capacidad de aceptar nuevas modalidades condicionales con una cantidad reducida de datos y recursos computacionales. Resuelve perfectamente la contradicción entre «no querer entrenar desde cero» y «no querer perder el conocimiento preentrenado», consolidándose como uno de los métodos más exitosos en el campo del ajuste fino de modelos difusivos.

\newpage
